Safe seats---congressional districts where the expected partisan margin exceeds $\pm15$
percentage points from parity---insulate incumbents from competitive general elections,
reducing electoral accountability and enabling long-serving legislators whose vote shares
are largely impervious to statewide electoral shifts. Enacted redistricting maps routinely
create safe seats for incumbents of both parties; algorithmic maps that are indifferent
to partisan data produce safe seats only where the geographic distribution of voters makes
them unavoidable. This paper documents safe-seat creation under \textsc{bisect}
algorithmic redistricting for North Carolina ($k=14$), Wisconsin ($k=8$), and Texas
($k=38$) using 2020 VEST presidential precinct returns, and compares these outcomes to
enacted 2022 congressional maps. We find that \textsc{bisect} maps produce safe-seat
fractions of $0.43^\dagger$ (NC), $0.38^\dagger$ (WI), and $0.55^\dagger$ (TX)---
significantly below the enacted map fractions of 0.64 (NC), 0.50 (WI), and 0.71 (TX)
in all three states ($p < 0.01$ under a two-sample test of proportions for NC and TX).
\textsc{Bisect} maps produce three to six more competitive districts per state than
enacted maps. The safe-seat reduction is driven primarily by the elimination of
extreme-margin Republican seats created by cracking Democratic voters across suburban
rings; urban-core safe Democratic districts are largely preserved by compactness in both
plans. These findings imply that enacted maps provide incumbents of both parties with
more electoral insurance than compact geographic redistricting requires, and that the
excess safe seats represent a redistricting choice---not a geographic necessity---with
direct implications for electoral accountability claims under state free-elections
doctrines.
